gin{equation}
\begin{aligned}
    \mathbf m^\star
    &=
    \frac{
        p\mathbf m_+
        -
        q\mathbf m_-
    }{
        p-q
    }
    \\
    &=
    \mathbf m_+
    +
    \frac{q}{p-q}
    \left(
        \mathbf m_+
        -
        \mathbf m_-
    \right)
    \\
    &=
    \mathbf m_+
    +
    \frac{\rho}{1-\rho}
    \left(
        \mathbf m_+
        -
        \mathbf m_-
    \right).
\end{aligned}
\label{eq:app_stable_extrapolation_identity}
\end{equation}
This proves Equation~\eqref{eq:signed_target_ratio}. In particular,
holding \(\mathbf m_+\) and \(\mathbf m_-\) fixed, the displacement from
the positive-only target is
\begin{equation}
    \mathbf m^\star-\mathbf m_+
    =
    \frac{\rho}{1-\rho}
    \left(
        \mathbf m_+
        -
        \mathbf m_-
    \right),
    \label{eq:app_equilibrium_displacement}
\end{equation}
whose magnitude increases with \(\rho\in(0,1)\) whenever
\(\mathbf m_+\neq\mathbf m_-\).

If \(\mathbf m^\star\notin\mathcal M\), no finite
\(\xi^\star\in\mathcal H\) can satisfy
Equation~\eqref{eq:app_equilibrium_condition}, and hence no finite
equilibrium exists.

More generally, consider any sequence
\(\mathbf m^\star_n\in\mathcal M\) that leaves every compact subset of
\(\mathcal M\), either by approaching its boundary or by becoming
unbounded, and define
\begin{equation}
    \xi^\star_n
    =
    (\nabla\psi)^{-1}
    \left(
        \mathbf m^\star_n
    \right).
\end{equation}
Suppose, for contradiction, that
\(\{\xi^\star_n\}\) remains in a compact subset
\(K\subset\mathcal H\). By continuity,
\(\nabla\psi(K)\) is a compact subset of \(\mathcal M\), but
\(\mathbf m^\star_n=\nabla\psi(\xi^\star_n)\) would then remain inside
that compact set. This contradicts the assumed behavior of
\(\mathbf m^\star_n\). Hence, the corresponding natural parameters must
leave every compact subset of \(\mathcal H\).

\medskip
\noindent\textbf{Critical regime.}
Suppose \(p=q>0\). Equation~\eqref{eq:app_aggregate_field} reduces to
\begin{equation}
\begin{aligned}
    \mathbf F(\xi)
    &=
    \bar{\mathbf m}_+
    -
    \bar{\mathbf m}_-
    \\
    &=
    p
    \left(
        \mathbf m_+
        -
        \mathbf m_-
    \right),
\end{aligned}
\label{eq:app_critical_field}
\end{equation}
which is independent of \(\xi\).

If \(\mathbf m_+\neq\mathbf m_-\), the continuous trajectory is
\begin{equation}
    \xi(t)
    =
    \xi(0)
    +
    tp
    \left(
        \mathbf m_+
        -
        \mathbf m_-
    \right),
    \label{eq:app_critical_continuous}
\end{equation}
and the discrete trajectory is
\begin{equation}
    \xi_t
    =
    \xi_0
    +
    t\alpha p
    \left(
        \mathbf m_+
        -
        \mathbf m_-
    \right).
    \label{eq:app_critical_discrete}
\end{equation}
Both exhibit persistent linear drift and admit no finite equilibrium.

If \(\mathbf m_+=\mathbf m_-\), then
\(\mathbf F(\xi)=\mathbf 0\) for every \(\xi\). Every point is
stationary, but a perturbed trajectory remains at the perturbed point
rather than returning to the original one. Hence, no point is an
isolated locally asymptotically stable equilibrium.

\medskip
\noindent\textbf{Negative-dominant regime.}
Suppose \(p<q\). Equation~\eqref{eq:app_signed_hessian} gives
\begin{equation}
    \nabla^2\mathcal L(\xi)
    =
    (q-p)
    \nabla^2\psi(\xi)
    \succ0.
\end{equation}
Thus, \(\mathcal L\) is strictly convex.

If a finite stationary point \(\xi^\star\) exists, the continuous field
Jacobian is
\begin{equation}
    J_{\mathbf F}(\xi^\star)
    =
    (q-p)
    \nabla^2\psi(\xi^\star)
    \succ0.
    \label{eq:app_negative_continuous_jacobian}
\end{equation}
All eigenvalues are positive, so the stationary point is repelling and
therefore unstable.

For the discrete update, the local Jacobian is
\begin{equation}
    J_{\mathrm{disc}}(\xi^\star)
    =
    I
    +
    \alpha(q-p)
    \nabla^2\psi(\xi^\star).
    \label{eq:app_negative_discrete_jacobian}
\end{equation}
Its eigenvalues are
\[
    1+\alpha(q-p)\lambda_i
    >
    1
\]
for every \(\alpha>0\) and every \(\lambda_i>0\). Therefore, every finite
stationary point is unstable under the discrete dynamics as well. If no
stationary point exists, there is trivially no finite stable equilibrium.

The four cases establish
Theorem~\ref{thm:aggregate_equilibria}.
\end{proof}


\subsection{Proof of
Theorem~\ref{thm:policy_family_manifestations}}
\label{app:policy_family_manifestations}

\begin{proof}
Let
\begin{equation}
    r
    =
    q-p
    >
    0.
\end{equation}

\medskip
\noindent\textbf{Gaussian policy.}
For a Gaussian policy with fixed covariance,
\begin{equation}
    \nabla_{\mu}
    \log\pi_{\mu}(a)
    =
    \Sigma^{-1}(a-\mu).
    \label{eq:app_gaussian_mean_score}
\end{equation}
The aggregate mean field is therefore
\begin{equation}
\begin{aligned}
    \mathbf{F}_{\mu}(\mu)
    &=
    \mathbb{E}_{\nu}
    \left[
        \widehat{A}(a)
        \Sigma^{-1}(a-\mu)
    \right] \\
    &=
    \Sigma^{-1}
    \left[
        p\mu_{+}
        -
        q\mu_{-}
        -
        (p-q)\mu
    \right].
\end{aligned}
\label{eq:app_gaussian_aggregate_field}
\end{equation}
Because $p-q=-r$ and
\begin{equation}
    \mu^{\dagger}
    =
    \frac{
        q\mu_{-}
        -
        p\mu_{+}
    }{
        r
    },
\end{equation}
Equation~\eqref{eq:app_gaussian_aggregate_field} becomes
\begin{equation}
    \mathbf{F}_{\mu}(\mu)
    =
    r\Sigma^{-1}
    \left(
        \mu-\mu^{\dagger}
    \right).
\end{equation}

For the continuous dynamics, define
\begin{equation}
    \delta(t)
    =
    \mu(t)-\mu^{\dagger}.
\end{equation}
Then
\begin{equation}
    \dot{\delta}(t)
    =
    r\Sigma^{-1}\delta(t),
\end{equation}
and hence
\begin{equation}
    \delta(t)
    =
    \exp
    \left(
        r\Sigma^{-1}t
    \right)
    \delta(0).
    \label{eq:app_gaussian_continuous_solution}
\end{equation}
Since $\Sigma^{-1}\succ0$,
\begin{equation}
    \left\|
        \delta(t)
    \right\|_2
    \geq
    \exp
    \left(
        r\lambda_{\min}
        \left(
            \Sigma^{-1}
        \right)t
    \right)
    \left\|
        \delta(0)
    \right\|_2.
    \label{eq:app_gaussian_continuous_growth}
\end{equation}
Therefore, unless
$\delta(0)=0$, the distance from the stationary point grows
exponentially and
$\|\mu(t)\|_2\rightarrow\infty$.

For the discrete update,
\begin{equation}
    \mu_{t+1}
    =
    \mu_t
    +
    \alpha
    \mathbf{F}_{\mu}(\mu_t),
\end{equation}
so
\begin{equation}
    \delta_{t+1}
    =
    \left(
        I
        +
        \alpha r\Sigma^{-1}
    \right)
    \delta_t.
\end{equation}
Iterating gives
\begin{equation}
    \delta_t
    =
    \left(
        I
        +
        \alpha r\Sigma^{-1}
    \right)^t
    \delta_0.
    \label{eq:app_gaussian_discrete_solution}
\end{equation}
All eigenvalues of the update matrix are strictly larger than one.
Consequently,
\begin{equation}
    \left\|
        \delta_t
    \right\|_2
    \geq
    \left[
        1
        +
        \alpha r
        \lambda_{\min}
        \left(
            \Sigma^{-1}
        \right)
    \right]^t
    \left\|
        \delta_0
    \right\|_2.
    \label{eq:app_gaussian_discrete_growth}
\end{equation}
Thus, every nonstationary discrete trajectory also diverges.

For any fixed historical action $a$,
\begin{equation}
\begin{aligned}
    \left\|
        \nabla_{\mu}
        \log\pi_{\mu}(a)
    \right\|_2
    &=
    \left\|
        \Sigma^{-1}(a-\mu)
    \right\|_2 \\
    &\geq
    \lambda_{\min}
    \left(
        \Sigma^{-1}
    \right)
    \left\|
        a-\mu
    \right\|_2.
\end{aligned}
\label{eq:app_gaussian_score_lower_bound}
\end{equation}
Since $\|\mu\|_2\rightarrow\infty$ while $a$ is fixed,
$\|a-\mu\|_2\rightarrow\infty$.
Equation~\eqref{eq:app_gaussian_score_lower_bound} therefore proves
Equation~\eqref{eq:gaussian_unbounded_score}.

\medskip
\noindent\textbf{Categorical policy.}
Let there be $K\geq2$ actions and write
\begin{equation}
    \pi_z
    =
    \operatorname{softmax}(z).
\end{equation}
Because adding a constant to every logit does not change the policy, we
work on the gauge-fixed subspace
\begin{equation}
    \mathcal{H}_0
    =
    \left\{
        z\in\mathbb{R}^{K}
        :
        \mathbf{1}^{\top}z=0
    \right\}.
    \label{eq:app_categorical_gauge}
\end{equation}

Let $\rho_{+}$ and $\rho_{-}$ denote the normalized positive and negative
weighted action distributions, and define
\begin{equation}
    b
    =
    p\rho_{+}
    -
    q\rho_{-}.
\end{equation}
The signed categorical objective can be written as
\begin{equation}
    \mathcal{L}(z)
    =
    b^{\top}z
    +
    r
    \log
    \left(
        \sum_{j=1}^{K}
        e^{z_j}
    \right)
    +
    C,
    \label{eq:app_categorical_objective}
\end{equation}
because $p-q=-r$. Its gradient and Hessian are
\begin{equation}
    \nabla\mathcal{L}(z)
    =
    b+r\pi_z
\end{equation}
and
\begin{equation}
    \nabla^2\mathcal{L}(z)
    =
    r
    \left[
        \operatorname{Diag}(\pi_z)
        -
        \pi_z\pi_z^{\top}
    \right].
    \label{eq:app_categorical_hessian}
\end{equation}
For every nonzero
$v\in\mathcal{H}_0$,
\begin{equation}
\begin{aligned}
    v^{\top}
    \nabla^2\mathcal{L}(z)
    v
    &=
    r
    \operatorname{Var}_{A\sim\pi_z}
    \left[
        v_A
    \right] \\
    &>
    0,
\end{aligned}
\end{equation}
because every softmax probability is positive and a nonzero vector in
$\mathcal{H}_0$ cannot be constant. Hence,
$\mathcal{L}$ is strictly convex on $\mathcal{H}_0$ and has at most one
finite stationary point.

Consider first the continuous gradient flow
\begin{equation}
    \dot{z}
    =
    \nabla\mathcal{L}(z).
\end{equation}
Along every trajectory,
\begin{equation}
    \frac{d}{dt}
    \mathcal{L}(z(t))
    =
    \left\|
        \nabla\mathcal{L}(z(t))
    \right\|_2^2
    \geq
    0.
    \label{eq:app_categorical_objective_growth}
\end{equation}
Suppose a nonstationary trajectory were bounded. Its closure would be
compact, so $\mathcal{L}$ would remain bounded above. Integrating
Equation~\eqref{eq:app_categorical_objective_growth} would then give
\begin{equation}
    \int_0^\infty
    \left\|
        \nabla\mathcal{L}(z(t))
    \right\|_2^2
    dt
    <
    \infty.
\end{equation}
Since the gradient is Lipschitz on the compact trajectory closure, the
standard gradient-flow argument implies
\begin{equation}
    \nabla\mathcal{L}(z(t))
    \longrightarrow
    0.
\end{equation}
Any accumulation point would therefore be a stationary point.

If no finite stationary point exists, this is an immediate
contradiction. If a stationary point $z^{\dagger}$ exists, strict
convexity gives, for every $z\neq z^{\dagger}$,
\begin{equation}
    (z-z^{\dagger})^{\top}
    \left[
        \nabla\mathcal{L}(z)
        -
        \nabla\mathcal{L}(z^{\dagger})
    \right]
    >
    0.
\end{equation}
Since
$\nabla\mathcal{L}(z^{\dagger})=0$,
\begin{equation}
    \frac{d}{dt}
    \frac{1}{2}
    \left\|
        z(t)-z^{\dagger}
    \right\|_2^2
    >
    0
\end{equation}
along every nonstationary trajectory. Such a trajectory cannot have a
subsequence converging to $z^{\dagger}$. Hence, the only bounded
continuous trajectory is the stationary one itself.

For the discrete update
\begin{equation}
    z_{t+1}
    =
    z_t
    +
    \alpha
    \nabla\mathcal{L}(z_t),
    \qquad
    \alpha>0,
\end{equation}
convexity gives
\begin{equation}
\begin{aligned}
    \mathcal{L}(z_{t+1})
    &\geq
    \mathcal{L}(z_t)
    +
    \nabla\mathcal{L}(z_t)^{\top}
    (z_{t+1}-z_t) \\
    &=
    \mathcal{L}(z_t)
    +
    \alpha
    \left\|
        \nabla\mathcal{L}(z_t)
    \right\|_2^2.
\end{aligned}
\label{eq:app_categorical_discrete_growth}
\end{equation}
If the trajectory were bounded,
Equation~\eqref{eq:app_categorical_discrete_growth} would imply
\begin{equation}
    \nabla\mathcal{L}(z_t)
    \longrightarrow
    0.
\end{equation}
A convergent subsequence would therefore approach a finite stationary
point. If no such point exists, this is impossible. If the unique
stationary point $z^{\dagger}$ exists, then for
$z_t\neq z^{\dagger}$,
\begin{equation}
\begin{aligned}
    \left\|
        z_{t+1}-z^{\dagger}
    \right\|_2^2
    &=
    \left\|
        z_t-z^{\dagger}
    \right\|_2^2 \\
    &\quad
    +
    2\alpha
    (z_t-z^{\dagger})^{\top}
    \nabla\mathcal{L}(z_t) \\
    &\quad
    +
    \alpha^2
    \left\|
        \nabla\mathcal{L}(z_t)
    \right\|_2^2 \\
    &>
    \left\|
        z_t-z^{\dagger}
    \right\|_2^2.
\end{aligned}
\end{equation}
The trajectory therefore cannot have a subsequence converging to
$z^{\dagger}$. Thus, every nonstationary discrete trajectory leaves
every compact subset of $\mathcal{H}_0$.

To relate logit divergence to the probability boundary, note that on
$\mathcal{H}_0$ the inverse softmax representation is
\begin{equation}
    z_i
    =
    \log\pi_z(i)
    -
    \frac{1}{K}
    \sum_{j=1}^{K}
    \log\pi_z(j).
    \label{eq:app_inverse_softmax}
\end{equation}
If the probabilities remained in a compact subset of the simplex
interior, all coordinates would be bounded below by some
$\varepsilon>0$, and
Equation~\eqref{eq:app_inverse_softmax} would imply bounded logits.
Therefore, an unbounded gauge-fixed logit trajectory must contain a
subsequence along which at least one action probability approaches zero.

Finally, for every action $a$,
\begin{equation}
    \nabla_z\log\pi_z(a)
    =
    e_a-\pi_z.
\end{equation}
Its squared norm satisfies
\begin{equation}
\begin{aligned}
    \left\|
        e_a-\pi_z
    \right\|_2^2
    &=
    \left(
        1-\pi_z(a)
    \right)^2
    +
    \sum_{j\neq a}
    \pi_z(j)^2 \\
    &\leq
    \left(
        1-\pi_z(a)
    \right)^2
    +
    \left(
        \sum_{j\neq a}
        \pi_z(j)
    \right)^2 \\
    &=
    2
    \left(
        1-\pi_z(a)
    \right)^2 \\
    &\leq
    2.
\end{aligned}
\end{equation}
This proves the categorical score bound and completes the proof.
\end{proof}

\subsection{Proof of Minimum-Order Control}
\label{app:min_order_control}

For fixed-covariance Gaussian mean coordinates,
Proposition~\ref{prop:distance_strength} gives
\(R(D)=\Theta(D-C_\Sigma)=\Theta(D)\) in the far field. Thus there are
constants \(0<c_1\le c_2<\infty\) and \(D_0\) such that, for
\(D\ge D_0\),
\[
    c_1 D
    \le
    R(D)
    \le
    c_2 D .
\]
If \(\omega(D)R(D)\) is bounded, then for some \(M<\infty\),
\[
    \omega(D)R(D)\le M
    \qquad
    \text{for all sufficiently large } D .
\]
Using the lower bound gives
\[
    \omega(D)
    \le
    \frac{M}{c_1D}
    =
    O(D^{-1}).
\]
Conversely, \(\omega(D)\le CD^{-1}\) and the upper bound imply
\[
    \omega(D)R(D)
    \le
    CD^{-1}c_2D
    =
    O(1).
\]
The same comparison with little-\(o\) proves that the weighted loop gain
vanishes exactly when \(\omega(D)=o(D^{-1})\). Finally,
\(D-C_\Sigma=r^2/2\) converts \(D^{-1}\) to the equivalent
reciprocal-quadratic order \(r^{-2}\).


\subsection{Empirical Negative-Update Utility}
\label{app:negative_update_utility}

To operationalize the utilit